% GENERATED FILE. Edit canonical lessons, then run npm run generate:books.
% canonical-source-hash: fnv1a64:0305dd75ea340ccf
% canonical-lessons: PT-C325-personagem, PT-C325-protagonista, PT-C325-heroi, PT-C325-vilao, PT-C325-desfecho

\chapter{Heróis e vilões}
\label{ch:pt-herois-e-viloes}
\hlchaptermodality{325}

\begin{chapteropening}
\textbf{By the end of this chapter:} \emph{I can talk about the characters of a story: the protagonist, the hero, the villain and the ending.}

Complete the last lesson of chapter 325: I can talk about the characters of a story: the protagonist, the hero, the villain and the ending.
\end{chapteropening}

\section[a personagem]{a personagem — the character (in a story)}

\label{lesson:PT-C325-personagem}



\begin{quote}
Before the new one: say the Portuguese for the narrator, then the Portuguese for the plot.
\end{quote}

\subsection*{You'll want to know: a personagem}
\textbf{a personagem} — \textquotedblleft{}the character (in a story)\textquotedblright{}.

\begin{quote}
A personagem principal é uma professora.
\end{quote}

\begin{quote}
\emph{The main character is a teacher.}
\end{quote}

In Portugal usually \emph{a personagem}; in Brazil also \emph{o personagem}.

\begin{cousinweb}
\emph{personagem} came from French \emph{personnage}, built on \emph{personne}, from Latin \emph{persōna} \textquotedblleft{}a mask, a role in a play\textquotedblright{} (Lewis \& Short).
\end{cousinweb}

\begin{grammarlens}[title={What you've built}]
One more word for telling, explaining and describing at B1.
\end{grammarlens}

\subsection*{Your turn}
\begin{itemize}
\raggedright
  \item \emph{Say it:} \emph{a personagem}
  \item \emph{Say it:} \emph{a personagem}, once more
  \item \emph{Say it:} a sentence of your own with \emph{a personagem}
  \item \emph{Recall:} say the Portuguese for the narrator, then the Portuguese for the plot, then say \emph{a personagem} again
\end{itemize}

\begin{tcolorbox}[breakable,colback=teal!4,colframe=teal!35!black,title={Before you move on}]
What does a personagem mean? (\textquotedblleft{}The character (in a story)\textquotedblright{}.) Say it once more.
\end{tcolorbox}
\section[o protagonista]{o protagonista — the protagonist, the main character}

\label{lesson:PT-C325-protagonista}



\begin{quote}
Before the new one: say the Portuguese for the plot, then the Portuguese for the character.
\end{quote}

\subsection*{You'll want to know: o protagonista}
\textbf{o protagonista} — \textquotedblleft{}the protagonist, the main character\textquotedblright{}.

\begin{quote}
O protagonista perde tudo no início da história.
\end{quote}

\begin{quote}
\emph{The protagonist loses everything at the start of the story.}
\end{quote}

The same form for a woman: \emph{a protagonista}.

\begin{cousinweb}
\emph{protagonista} goes back to Greek \emph{prōtagōnist\'{\={e}}s} \textquotedblleft{}the first actor\textquotedblright{}, from \emph{prôtos} \textquotedblleft{}first\textquotedblright{} and \emph{agōnist\'{\={e}}s} \textquotedblleft{}a contestant\textquotedblright{} (EtymDB has no entry).
\end{cousinweb}

\begin{grammarlens}[title={What you've built}]
One more word for telling, explaining and describing at B1.
\end{grammarlens}

\subsection*{Your turn}
\begin{itemize}
\raggedright
  \item \emph{Say it:} \emph{o protagonista}
  \item \emph{Say it:} \emph{o protagonista}, once more
  \item \emph{Say it:} a sentence of your own with \emph{o protagonista}
  \item \emph{Recall:} say the Portuguese for the plot, then the Portuguese for the character, then say \emph{o protagonista} again
\end{itemize}

\begin{tcolorbox}[breakable,colback=teal!4,colframe=teal!35!black,title={Before you move on}]
What does o protagonista mean? (\textquotedblleft{}The protagonist, the main character\textquotedblright{}.) Say it once more.
\end{tcolorbox}
\section[o herói]{o herói — the hero}

\label{lesson:PT-C325-heroi}



\begin{quote}
Before the new one: say the Portuguese for the character, then the Portuguese for the protagonist.
\end{quote}

\subsection*{You'll want to know: o herói}
\textbf{o herói} — \textquotedblleft{}the hero\textquotedblright{}.

\begin{quote}
No fim, o herói salva a cidade.
\end{quote}

\begin{quote}
\emph{In the end the hero saves the city.}
\end{quote}

A woman is \emph{a heroína}.

\begin{cousinweb}
\emph{herói} is a learned borrowing from Latin \emph{hērōs} (EtymDB), from Greek \emph{h\'{\={e}}rōs} \textquotedblleft{}a demigod, a hero\textquotedblright{}.
\end{cousinweb}

\begin{grammarlens}[title={What you've built}]
One more word for telling, explaining and describing at B1.
\end{grammarlens}

\subsection*{Your turn}
\begin{itemize}
\raggedright
  \item \emph{Say it:} \emph{o herói}
  \item \emph{Say it:} \emph{o herói}, once more
  \item \emph{Say it:} a sentence of your own with \emph{o herói}
  \item \emph{Recall:} say the Portuguese for the character, then the Portuguese for the protagonist, then say \emph{o herói} again
\end{itemize}

\begin{tcolorbox}[breakable,colback=teal!4,colframe=teal!35!black,title={Before you move on}]
What does o herói mean? (\textquotedblleft{}The hero\textquotedblright{}.) Say it once more.
\end{tcolorbox}
\section[o vilão]{o vilão — the villain}

\label{lesson:PT-C325-vilao}



\begin{quote}
Before the new one: say the Portuguese for the protagonist, then the Portuguese for the hero.
\end{quote}

\subsection*{You'll want to know: o vilão}
\textbf{o vilão} — \textquotedblleft{}the villain\textquotedblright{}.

\begin{quote}
O vilão do filme é muito engraçado.
\end{quote}

\begin{quote}
\emph{The villain in the film is very funny.}
\end{quote}

A woman is \emph{a vilã}.

\begin{cousinweb}
\emph{vilão} comes from Late Latin \emph{vīllānus} \textquotedblleft{}a farm worker\textquotedblright{} (EtymDB), from \emph{vīlla} \textquotedblleft{}a farm\textquotedblright{}: a peasant, later a scoundrel.
\end{cousinweb}

\begin{grammarlens}[title={What you've built}]
One more word for telling, explaining and describing at B1.
\end{grammarlens}

\subsection*{Your turn}
\begin{itemize}
\raggedright
  \item \emph{Say it:} \emph{o vilão}
  \item \emph{Say it:} \emph{o vilão}, once more
  \item \emph{Say it:} a sentence of your own with \emph{o vilão}
  \item \emph{Recall:} say the Portuguese for the protagonist, then the Portuguese for the hero, then say \emph{o vilão} again
\end{itemize}

\begin{tcolorbox}[breakable,colback=teal!4,colframe=teal!35!black,title={Before you move on}]
What does o vilão mean? (\textquotedblleft{}The villain\textquotedblright{}.) Say it once more.
\end{tcolorbox}
\section[o desfecho]{o desfecho — the ending, the outcome}

\label{lesson:PT-C325-desfecho}



\begin{quote}
Before the new one: say the Portuguese for the hero, then the Portuguese for the villain.
\end{quote}

\subsection*{You'll want to know: o desfecho}
\textbf{o desfecho} — \textquotedblleft{}the ending, the outcome\textquotedblright{}.

\begin{quote}
O desfecho da história surpreendeu toda a gente.
\end{quote}

\begin{quote}
\emph{The ending of the story surprised everyone.}
\end{quote}

\emph{Um desfecho feliz} — a happy ending.

\begin{cousinweb}
\emph{desfecho} is built on \emph{desfechar}, from \emph{des-} and \emph{fechar} \textquotedblleft{}to close\textquotedblright{}: the knot of the story comes undone. EtymDB traces \emph{fechar} through Vulgar Latin \emph{fesclāre} to Latin \emph{fistulāre}, built on \emph{fistula} \textquotedblleft{}a pipe, a tube\textquotedblright{} (Lewis \& Short).
\end{cousinweb}

\begin{grammarlens}[title={What you've built}]
One more word for telling, explaining and describing at B1.
\end{grammarlens}

\subsection*{Your turn}
\begin{itemize}
\raggedright
  \item \emph{Say it:} \emph{o desfecho}
  \item \emph{Say it:} \emph{o desfecho}, once more
  \item \emph{Say it:} a sentence of your own with \emph{o desfecho}
  \item \emph{Recall:} say the Portuguese for the hero, then the Portuguese for the villain, then say \emph{o desfecho} again
\end{itemize}

\begin{tcolorbox}[breakable,colback=teal!4,colframe=teal!35!black,title={Before you move on}]
What does o desfecho mean? (\textquotedblleft{}The ending, the outcome\textquotedblright{}.) Say it once more.
\end{tcolorbox}
